\section{De 3-Gram a GPT-3: la escala de calidad de los modelos de lenguaje}

Un \term{language model} (modelo de lenguaje) asigna probabilidades a secuencias de tokens. Dado $x_1,\ldots,x_T$, la factorización autorregresiva es
\[
p(x_1,\ldots,x_T)=\prod_{t=1}^{T}p(x_t\mid x_{<t}).
\]
Las diferencias entre modelos están principalmente en cómo se representa el contexto condicional, cuánta historia pueden usar y cómo escalan los parámetros y los datos. El ponente muestra primero samples, y no primero benchmarks, para que el lector vea de manera intuitiva los modos de falla de los modelos probabilísticos.

\subsection{3-Gram: la estadística local genera «palabras relacionadas», pero difícilmente forma una semántica global}

Un \term{n-gram} supone que el token actual depende solo de los $n-1$ tokens anteriores. Un 3-gram se escribe $p(x_t\mid x_{t-2},x_{t-1})$. Sus probabilidades pueden obtenerse suavizando conteos, y es fácil de entrenar e interpretar; sin embargo, su ventana de contexto es fija, las combinaciones raras son dispersas y no puede mantener entidades, hechos ni estructura discursiva a larga distancia.

\lecturefigure{slide-01-3gram.jpg}{El sample de 3-gram ya tiene palabras localmente relacionadas, pero carece de sintaxis estable y de semántica global.}{Video oficial 00:08; encuadre histórico de Shannon 1951.}

\begin{knowledgebox}{Lectura de la figura: no basta con preguntar si la oración «parece inglés»}
En el sample, montos, años y vocabulario comercial coocurren localmente, lo que indica que la estadística de ventana corta captura la adyacencia de palabras; en cambio, las relaciones numéricas, la consistencia de entidades y el propósito de la oración se derrumban. La evaluación puede dividirse en cuatro niveles: lexical coherence, syntactic coherence, semantic consistency y factual consistency. Un n-gram suele tener señal en el corto alcance de los dos primeros niveles; en los dos últimos le falta estado.
\end{knowledgebox}

\begin{knowledgebox}{Definición y limitaciones de la perplexity}
La \term{perplexity} (perplejidad) es la exponencial de la cross-entropy promedio:
\[
\operatorname{PPL}=\exp\!\left(-\frac{1}{T}\sum_{t=1}^{T}\log p(x_t\mid x_{<t})\right).
\]
Intuitivamente puede entenderse como el «número efectivo de candidatos» en cada paso; un valor más bajo suele indicar mejor predicción sobre esa distribución de prueba. Sin embargo, no conviene comparar directamente la PPL entre distintos tokenizers, limpiezas de datos y dominios; una PPL más baja tampoco garantiza por sí misma la exactitud factual, la controlabilidad ni el éxito en la tarea.
\end{knowledgebox}

\subsection{RNN: un estado oculto en lugar de una ventana fija}

Esta sección pasa de la ventana fija al estado recurrente. La RNN actualiza $h_t=f(x_t,h_{t-1})$, de modo que, en teoría, cualquier historia puede influir en la predicción actual. Elimina la ventana fija del n-gram, pero comprime toda la historia en un único estado y obliga a que el entrenamiento sea secuencial en el tiempo. Las secuencias largas además enfrentan desvanecimiento o explosión del gradiente y sobrescritura de información.

\lecturefigure{slide-02-rnn.jpg}{El sample de la RNN es más coherente que el de 3-gram, pero las entidades y la estructura de largo alcance siguen derivando con facilidad.}{Video oficial 00:27; Sutskever et al. 2011.}

\begin{knowledgebox}{Lectura de la figura: lo que mejora la RNN es la continuidad del estado}
La RNN puede sostener temas locales y patrones de caracteres o palabras durante más tiempo, por lo que el sample ya no parece solo una bolsa de palabras. El problema es que cada paso pasa por el mismo recurrent bottleneck, y la información antigua debe atravesar muchas transformaciones no lineales para llegar al futuro. Que en la salida aparezcan «oraciones que parecen lenguaje y argumentos que no parecen argumentos» muestra que la fluidez y la planeación global siguen siendo capacidades distintas.
\end{knowledgebox}

\teachervoice{El ponente muestra los modelos nivel por nivel con samples reales, en lugar de afirmar que «la nueva arquitectura supera en todo a la anterior». Nota de clase: el progreso de los modelos generativos suele manifestarse primero en que los errores se vuelven más difíciles de detectar, por lo que la evaluation también debe actualizarse a la par.}

\subsection{Big LSTM: las compuertas y la escala mejoran la memoria, pero no eliminan la ruta secuencial}

La \term{LSTM} (Long Short-Term Memory) usa input, forget y output gates para controlar el cell state, lo que mitiga los problemas de optimización de la RNN simple. LSTM más grandes, más datos y better regularization impulsaron de forma notable el language modeling, pero el entrenamiento sigue siendo recursivo por pasos de tiempo, y la información de larga distancia todavía debe recorrer una ruta de $O(T)$.

\lecturefigure{slide-03-big-lstm.jpg}{La LSTM a gran escala genera texto más parecido al natural, pero aún exhibe repeticiones, deriva y errores estructurales.}{Video oficial 01:12; Jozefowicz et al. 2016.}

\begin{knowledgebox}{Lectura de la figura: scale y architecture cambian el sample al mismo tiempo}
La slide coloca a la Big LSTM entre la RNN y el Transformer, como recordatorio de no atribuir todo el progreso a un solo componente. Las compuertas mejoran la memoria, la escala del modelo y de los datos reduce la incertidumbre local y la training recipe mejora la estabilidad; pero la recurrent dependency sigue limitando el paralelismo en hardware. Una comparación justa exige controlar el número de parámetros, el número de tokens y el compute, en lugar de elegir solo las muestras más vistosas.
\end{knowledgebox}

\subsection{Transformer: entrenamiento paralelo y rutas de información más cortas}

El Transformer usa self-attention para que los tokens interactúen directamente; las operaciones matriciales se adaptan mejor a los accelerators y permiten escalar a lotes, modelos y datos más grandes. No resolvió automáticamente los problemas de hechos, planeación y objetivos, pero cambió el scaling regime factible.

\lecturefigure{slide-04-transformer.jpg}{La mejora del sample del Transformer proviene de la acción conjunta de arquitectura, escala, datos y entrenamiento.}{Video oficial 01:52.}

\begin{knowledgebox}{Lectura de la figura: la model completion es una distribución condicional, no una consulta a una base de conocimiento}
La slide presenta un prompt y una completion: la salida puede continuar el formato y el tema, pero aun así inventar detalles. Un language model predice «qué continuación es probable dentro de la distribución de entrenamiento» y no ejecuta un fact-check por defecto. A medida que los samples se vuelven más fluidos, a las personas les resulta más fácil confundir calidad lingüística con calidad factual; por eso la later evaluation requiere grounding en la tarea, retrieval o external checks.
\end{knowledgebox}

\subsection{GPT-2: Transformer escalado y muestras de texto largo}

Esta sección pasa al decoder-only scaling. GPT-2 escala el Transformer a más datos y parámetros, y pone el énfasis en el zero-shot task behavior. En clase se muestra primero la escala del modelo y de los datos, y después muestras de noticias, con la intención de que el lector observe la coherence entre párrafos, la imitación de estilo y las failures.

\lecturefigure{slide-05-gpt2-model.jpg}{GPT-2 amplía el Transformer decoder-only hasta convertirlo en un modelo general de generación de texto.}{Video oficial 02:33; Radford et al. 2019.}

\begin{knowledgebox}{Lectura de la figura: el número de parámetros es solo una columna del sistema}
Aumentar la escala del modelo suele ir acompañado de más tokens de entrenamiento, distintos lote/optimizer, context y compute. Decir únicamente «tiene más parámetros, por eso aparece la capacidad» ignora la data diversity y la optimization. Al leer esta slide, conviene tomar architecture, parameters, data, objective y evaluation como una recipe conjunta.
\end{knowledgebox}

\lecturefigure{slide-06-gpt2-example-1.jpg}{La muestra de noticias de GPT-2 exhibe estilo discursivo y continuidad de entidades, y al mismo tiempo conserva un riesgo de invención verificable.}{Video oficial 03:05.}

\begin{knowledgebox}{Lectura de la figura: las muestras largas se revisan claim por claim}
Primero se revisan la sintaxis y la conexión entre párrafos; después se marcan las claims sobre personas, lugares, cifras y causas; por último se verifica si esas claims pueden sustentarse con el prompt o con evidencia externa. El modelo puede mantener el estilo periodístico y, aun así, «completar con seguridad» los detalles. Si la evaluación de muestras solo pregunta si el texto se lee con fluidez, pasará por alto de manera sistemática la hallucination.
\end{knowledgebox}

\lecturefigure{slide-07-gpt2-example-2.jpg}{La segunda muestra de GPT-2 muestra que distintos prompts exponen distintos tipos de coherence failure.}{Video oficial 03:38.}

\begin{knowledgebox}{Lectura de la figura: no se debe representar la distribución del modelo con una sola muestra}
El sampling seed, la temperature, el top-p y el prompt modifican la salida. Una muestra vistosa no representa el desempeño promedio, y una muestra mala tampoco prueba que el modelo sea inútil. Es necesario fijar el decoding protocol, reportar el número de muestras y combinar likelihood, human rating y downstream success.
\end{knowledgebox}

\subsection{GPT-3: ¿pueden las personas identificar noticias generadas?}

Cuando la calidad del texto mejora, la evaluation pasa de «se ve claramente mal» a experimentos controlados con personas. El artículo de GPT-3 pidió a los participantes juzgar si una noticia había sido generada por el modelo; la precisión varía con la escala del modelo. Este experimento mide la detection bajo condiciones específicas de prompt, longitud, muestreo y participantes; no equivale a veracidad ni a inteligencia general.

\lecturefigure{slide-08-gpt3-human-detection.jpg}{GPT-3 convierte el sample realism en un experimento cuantificable mediante una tarea de identificación por personas.}{Video oficial 05:12; Brown et al. 2020.}

\begin{knowledgebox}{Lectura de la figura: primero se define la unidad experimental}
Cada item es un fragmento de noticia generado por el modelo o escrito por una persona; el participante emite un juicio de clasificación binaria. El resultado depende de la familiaridad con el tema, la longitud del texto, la temperatura de generación y la instruction. Un valor cercano a chance solo indica que, en ese contexto, es difícil distinguir el origen; no indica que el contenido sea verídico, inofensivo o apto para la producción de noticias.
\end{knowledgebox}

\lecturefigure{slide-09-gpt3-detection-curve.jpg}{Los modelos más grandes dificultan la detección, pero la curva debe interpretarse junto con el diseño experimental.}{Video oficial 08:10; Brown et al. 2020.}

\begin{knowledgebox}{Lectura de la figura: el eje horizontal es la escala; el vertical no es «inteligencia»}
La curva relaciona el model size con la human detection accuracy. La tendencia principal es que, al aumentar la escala, el estilo generado se acerca más a la distribución humana; la curva no mide la exactitud factual, la trazabilidad de las fuentes ni la consistencia a largo plazo. Las barras de error, el baseline y la sample policy determinan la solidez de la conclusión. No se puede extrapolar una sola curva a todos los idiomas y a todas las tareas de generación.
\end{knowledgebox}

\subsection{Resumen de la sección}

Del n-gram a GPT-3, los samples se vuelven progresivamente más coherentes, y la evaluación pasa de fallas evidentes a simple vista a experimentos controlados con personas. La scale cambia la calidad de la distribución, pero hace aún más importantes los problemas de veracidad y de identificación del origen.
